% ---------------------------------------------------------------
% Map figures rendered from the interactive Leaflet maps
% (regenerate with:  python render_maps_to_png.py)
%
% Preamble needs:  \usepackage{graphicx}
%                  \graphicspath{{figures_maps/}}   % or copy PNGs next to the .tex
% ---------------------------------------------------------------

% === Section 5.2 -- "Figure?? maps the resulting locational value dV ..." =========
\begin{figure}[htbp]
  \centering
  \includegraphics[width=0.78\textwidth]{map_siting_score_clean.png}
  \caption[Locational value of a reference battery]{%
    Locational value $\Delta V$ of a reference 50\,MW/200\,MWh battery across Germany,
    derived from geolocated redispatch events (2013--2025, confidence $>0.5$).
    Colour: annual $\Delta V$ per candidate grid cell in million~EUR/year;
    black circles mark the ten highest-scoring cells.
    Red/blue lines show the 380/220\,kV transmission grid.}
  \label{fig:siting_score}
\end{figure}

% === Section 5.3 -- "the 17 German PyPSA nodes (circles in Figure??)" ============
\begin{figure}[htbp]
  \centering
  \includegraphics[width=0.78\textwidth]{map_siting_vs_pypsa_clean.png}
  \caption[Siting score vs.\ PyPSA-Eur nodal benchmark]{%
    Redispatch-based siting score (background heatmap) against the PyPSA-Eur nodal
    arbitrage benchmark (circles, one per German node; size and colour scale with the
    node's arbitrage value). The two rankings agree with a Spearman correlation of
    $\rho = 0.65$ ($\tau = 0.43$, $p = 0.005$).}
  \label{fig:siting_vs_pypsa}
\end{figure}

% === Section 5.1 -- redispatch events / grid stress backdrop ====================
\begin{figure}[htbp]
  \centering
  \includegraphics[width=0.78\textwidth]{map_redispatch_kde_clean.png}
  \caption[Spatial distribution of redispatch events]{%
    Kernel-density estimate of geolocated redispatch events in Germany.
    Only TSO (\"UNB) redispatch is covered; DSO-level measures are not included.}
  \label{fig:redispatch_kde}
\end{figure}

% === Section 5.2 -- optional: score vs. where batteries actually are =============
\begin{figure}[htbp]
  \centering
  \includegraphics[width=0.78\textwidth]{map_siting_vs_batteries_clean.png}
  \caption[Siting score vs.\ realised battery locations]{%
    Siting score (background) against actual MaStR battery locations
    (green: in operation, orange: planned; marker size $\propto \sqrt{\text{MW}}$).
    Realised siting follows land availability and connection queues rather than
    locational grid value.}
  \label{fig:siting_vs_batteries}
\end{figure}

% === Appendix A/B -- supplementary =============================================
\begin{figure}[htbp]
  \centering
  \includegraphics[width=0.78\textwidth]{map_source_nodes_clean.png}
  \caption[Geolocated redispatch source nodes]{%
    Geolocated redispatch source nodes, coloured by net direction
    (net upward vs.\ net downward dispatch).}
  \label{fig:source_nodes}
\end{figure}

\begin{figure}[htbp]
  \centering
  \includegraphics[width=0.78\textwidth]{map_h_sweep_clean.png}
  \caption[Siting score at the default kernel bandwidth]{%
    Siting score at the default kernel bandwidth $h$ (one frame of the bandwidth sweep,
    see Appendix~\ref{app:bandwidth}).}
  \label{fig:h_sweep_map}
\end{figure}

% --- Two maps side by side (alternative layout) --------------------------------
% \begin{figure}[htbp]
%   \centering
%   \begin{subfigure}[t]{0.48\textwidth}
%     \includegraphics[width=\linewidth]{map_siting_score_clean.png}
%     \caption{Redispatch-based siting score}
%   \end{subfigure}\hfill
%   \begin{subfigure}[t]{0.48\textwidth}
%     \includegraphics[width=\linewidth]{map_siting_vs_pypsa_clean.png}
%     \caption{PyPSA-Eur nodal benchmark overlay}
%   \end{subfigure}
%   \caption{Siting score and benchmark.}
% \end{figure}
% (requires \usepackage{subcaption})
